\documentclass[10pt,twocolumn]{article}
\usepackage[T1]{fontenc}
\usepackage[utf8]{inputenc}
\usepackage{lmodern}
\usepackage[margin=1.8cm,top=2.4cm,bottom=2cm,columnsep=0.6cm]{geometry}
\usepackage{fancyhdr}
\usepackage{titlesec}
\usepackage{booktabs}
\usepackage{amsmath,amssymb,amsfonts}
\usepackage{microtype}
\usepackage{xcolor}
\usepackage{mdframed}
\usepackage{parskip}
\usepackage{enumitem}
\usepackage{hyperref}

\definecolor{rulered}{RGB}{180,30,30}
\definecolor{boxgray}{RGB}{245,245,245}
\definecolor{keywordgray}{RGB}{100,100,100}

\pagestyle{fancy}
\fancyhf{}
\fancyhead[L]{\textsc{\small Claude Mag}}
\fancyhead[C]{\small\textit{Machine Learning Research Digest}}
\fancyhead[R]{\small 2026-09-28}
\fancyfoot[C]{\small\thepage}
\renewcommand{\headrulewidth}{0.8pt}

\titleformat{\section}{\normalsize\bfseries\color{rulered}}{}{0pt}{}%
  [\vspace{-4pt}\color{rulered}\rule{\columnwidth}{0.4pt}]
\titlespacing{\section}{0pt}{8pt}{4pt}

\hypersetup{colorlinks=true,urlcolor=rulered,linkcolor=black}

\begin{document}
\setlength{\parindent}{0pt}
\setlength{\parskip}{4pt}

{\small\color{keywordgray}\textbf{Keywords:}
  scaling laws \textbullet{}
  inference-time compute \textbullet{}
  chain-of-thought \textbullet{}
  Bayesian analysis}\par\vspace{4pt}

{\LARGE\bfseries\raggedright
  Scaling of Capability and Efficiency}\par\vspace{4pt}

{\large\itshape\raggedright
  Joint Quantitative Laws for Solve Rate and Token Cost
  in Large Reasoning Models}\par\vspace{6pt}

{\small\textbf{By} (authors not specified in pre-print metadata)}\par\vspace{2pt}

{\small\color{keywordgray}arXiv:2609.27166 \textbullet{} September 2026}
\par\vspace{2pt}\noindent{\color{rulered}\rule{\columnwidth}{0.6pt}}\vspace{6pt}

Using hierarchical Bayesian models applied to the DeepSeek-R1-Distill family,
this paper shows that solve probability decays exponentially with problem
hardness while the characteristic scale grows as a power law with model
parameters---and that among solved instances, output length grows as a power
law with instance size, providing the first joint quantitative account of
capability and efficiency scaling at inference time.

\begin{mdframed}[backgroundcolor=boxgray,linecolor=rulered,linewidth=1pt,
    innerleftmargin=6pt,innerrightmargin=6pt,
    innertopmargin=6pt,innerbottommargin=6pt]
\textbf{\small AT A GLANCE}\par\vspace{4pt}
\begin{description}[leftmargin=0pt,labelsep=4pt,itemsep=2pt,parsep=0pt,
    font=\small\bfseries]
  \item[Problem:] Capability and efficiency in chain-of-thought reasoning
    both depend on problem hardness and model size, but their joint scaling
    behaviour has not been quantitatively characterized.
  \item[Why it matters:] Without joint scaling laws, predicting compute
    requirements for a given problem difficulty and model size is guesswork.
  \item[Method:] Hierarchical Bayesian models fit to the DeepSeek-R1-Distill
    family across four algorithmic problem classes with controlled hardness.
  \item[Result:] Exponential capability decay with power-law characteristic
    scale; power-law token cost growth among correctly solved instances.
  \item[Evidence:] 1.5B--70B parameter DeepSeek-R1-Distill models on
    arithmetic, multiplication, sorting, and path-finding.
\end{description}
\end{mdframed}

\section{The Big Picture}

Large reasoning models (LRMs) achieve state-of-the-art results by generating
extended chain-of-thought traces before producing final answers. The amount
of computation consumed at inference time---measured in output tokens---grows
substantially with problem difficulty, and harder problems are solved by
larger models. Both capability (solve rate) and efficiency (tokens per
correct solution) thus depend on at least two variables: problem hardness and
model size.

Despite the obvious importance of understanding this dependency, it has not
been characterized in the form of quantitative scaling laws. Practitioners
must therefore guess how many tokens a given model will require to solve
problems of a given difficulty, and how much larger a model must be to solve
problems that the current model cannot.

This paper provides the first joint quantitative account of capability and
efficiency scaling at inference time, using controlled algorithmic problem
classes where hardness is operationalized as instance size and correctness
is verified exactly.

\section{Why Existing Approaches Fall Short}

\textbf{Average accuracy evaluations} aggregate performance across problem
difficulties, masking the exponential decay with hardness and conflating
easy-instance dominance with true capability on hard instances.

\textbf{Single-model comparisons} at fixed compute budgets cannot reveal
how model scale shifts the capability-efficiency frontier.

\textbf{Simple regression on token counts} ignores the joint structure of
success and token cost, treating failed attempts identically to efficient
successes.

\textbf{Prior scaling laws} (Chinchilla, etc.) address training compute and
not inference-time token generation conditioned on problem difficulty, leaving
the test-time compute dimension uncharacterized.

\section{Core Mechanism}

\textbf{Hierarchical Bayesian Model.} The paper fits a generative model
to solve indicators and token counts across models and problem instances.
The model has three layers:
\begin{enumerate}[itemsep=2pt,parsep=0pt]
  \item \textbf{Instance level:} binary solve indicator and token count per
    instance per model.
  \item \textbf{Class level:} partial pooling of scaling exponents across
    four problem classes.
  \item \textbf{Model level:} power-law relationship between model size and
    scaling parameters.
\end{enumerate}

\textbf{Problem Classes.} Four algorithmic tasks with exact verifiers:
arithmetic computation, digit multiplication, list sorting, and path-finding.
Hardness is controlled by instance size $n$ (digits, list length, or graph
nodes).

\textbf{Model Family.} DeepSeek-R1-Distill at 1.5B, 7B, 14B, 32B, and 70B
parameters---a controlled family where differences arise from model scale
alone.

\section{Technical Formulation}

Solve probability at instance size $n$ and model parameter count $m$:
\[
p(n, m) = p_0(m) \cdot \exp\!\left(-\frac{n}{\lambda(m)}\right)
\]

Characteristic scale as a power law in model parameters:
\[
\lambda(m) = \alpha \cdot m^\beta
\]

Conditional output length (correctly solved instances):
\[
\mathbb{E}[L \mid \text{solved},\, n,\, m]
= \gamma(m) \cdot n^{\delta}
\]

Hierarchical priors (partial pooling across problem classes $k=1,\ldots,K$):
\[
\beta_k \sim \mathcal{N}(\mu_\beta,\, \sigma_\beta^2),
\quad
\delta_k \sim \mathcal{N}(\mu_\delta,\, \sigma_\delta^2)
\]

All parameters are jointly estimated by MCMC, yielding posterior distributions
over scaling exponents $\beta$ and $\delta$ per class.

\section{Learning or Inference Procedure}

\begin{enumerate}[itemsep=2pt,parsep=0pt]
  \item For each (model, problem class, instance size): generate up to $R$
    attempts; record solve indicator and output token count.
  \item Fit the hierarchical Bayesian model to the full dataset using MCMC
    (Stan or equivalent).
  \item Report posterior means and 90\% credible intervals for
    $\alpha, \beta, \gamma, \delta$ per class.
  \item Use fitted model to predict solve probability and expected token cost
    for new (model size, instance size) pairs.
\end{enumerate}

\section{What the Guarantee Says}

The hierarchical model provides posterior uncertainty over scaling exponents.
Within the linear-arithmetic and combinatorial-algorithmic problem classes
studied, exponential solve-rate decay with power-law characteristic scale is
a consistent structural finding, not an artifact of a single model or
problem class. Cross-class partial pooling provides regularization that
stabilizes estimates at large difficulty where data is sparse.

\section{Experimental Findings}

\begin{itemize}[itemsep=2pt,parsep=0pt]
  \item \textbf{5 models:} 1.5B to 70B DeepSeek-R1-Distill
  \item \textbf{4 problem classes:} arithmetic, digit multiplication, list
    sorting, path-finding
  \item Solve probability fits exponential decay with hardness ($R^2 > 0.92$)
    across all classes
  \item Characteristic scale exponent $\beta \approx 0.3$--$0.5$ by class;
    path-finding has the steepest decay
  \item DeepSeek-R1-Distill-70B achieves approximately $10\times$ larger
    characteristic scale than the 1.5B model on arithmetic
  \item Output length exponent $\delta \approx 1.0$--$1.5$: near-linear to
    super-linear growth with instance size
  \item Efficiency degrades faster than capability with hardness: diminishing
    returns to scale on the hardest instances
\end{itemize}

\section{Ablations and Interpretation}

\textbf{Problem class differences:} Path-finding shows the steepest capability
decay, suggesting it requires qualitatively more extensive search. Arithmetic
shows the shallowest, consistent with near-algorithmic solutions.

\textbf{Pooling vs.\ no pooling:} Hierarchical pooling across classes stabilizes
estimates at large difficulty; without pooling, estimates at the hardest
difficulty levels have large variance.

\textbf{Efficiency vs.\ capability trade-off:} The power-law exponent
$\delta > 1$ for output length implies that even when a larger model can solve
harder instances, it pays a super-linear token cost. Compute-optimal scaling
must balance capability and efficiency gains.

\textbf{Implications:} Given model size $m$ and a target solve probability
$p^*$ on instances of size $n^*$, the minimum required model size is:
$m^* = (\lambda^{-1}(-n^*/\ln p^*)/\alpha)^{1/\beta}$---a prediction the
framework makes interpretable and calibrated.

\section*{Reference}

\textbf{Scaling of Capability and Efficiency at Inference Time in Large
Reasoning Models}. arXiv:2609.27166, September 2026.\\
\url{https://arxiv.org/abs/2609.27166}

\end{document}
